\subsection{Optimization Objective Comparison: MSE vs. Correlation}
\label{sec:mse_vs_pearson}

A critical design choice in learning mergeability prediction coefficients is the optimization objective. We compare two objectives: maximizing Pearson correlation versus minimizing mean squared error (MSE).

\paragraph{Objective Functions.}
For Pearson correlation maximization, we optimize:
\begin{equation}
    \mathcal{L}_{\text{Pearson}} = -\frac{\sum_i (y_i - \bar{y})(\hat{y}_i - \bar{\hat{y}})}{\sqrt{\sum_i (y_i - \bar{y})^2 \sum_i (\hat{y}_i - \bar{\hat{y}})^2}} + \lambda \|\mathbf{w}\|_1
\end{equation}
where $y_i$ is the actual merge performance, $\hat{y}_i = \mathbf{w}^\top \mathbf{x}_i$ is the predicted mergeability, and $\mathbf{w}$ are the learned coefficients.

For MSE minimization:
\begin{equation}
    \mathcal{L}_{\text{MSE}} = \frac{1}{N}\sum_i (y_i - \hat{y}_i)^2
\end{equation}

\paragraph{Results.}
Tables~\ref{tab:objective_comparison_b16} and~\ref{tab:objective_comparison_b32} present the validation Pearson correlation for each objective across all merging methods.

\begin{table}[h]
\centering
\caption{ViT-B-16: Validation Pearson correlation ($r$) for different optimization objectives. Higher is better.}
\label{tab:objective_comparison_b16}
\begin{tabular}{@{}lccc@{}}
\toprule
\textbf{Method} & \textbf{Pearson + L1} & \textbf{MSE} & \textbf{Difference} \\
\midrule
Weight Avg & 0.589 & 0.244 & $+0.345$ \\
Arithmetic & 0.445 & 0.061 & $+0.384$ \\
TSV & 0.627 & 0.187 & $+0.440$ \\
TIES & 0.463 & 0.098 & $+0.365$ \\
DARE & 0.596 & 0.143 & $+0.453$ \\
\midrule
\textbf{Average} & \textbf{0.544} & 0.147 & $+0.397$ \\
\bottomrule
\end{tabular}
\end{table}

\begin{table}[h]
\centering
\caption{ViT-B-32: Validation Pearson correlation ($r$) for different optimization objectives. Higher is better.}
\label{tab:objective_comparison_b32}
\begin{tabular}{@{}lccc@{}}
\toprule
\textbf{Method} & \textbf{Pearson + L1} & \textbf{MSE} & \textbf{Difference} \\
\midrule
Weight Avg & 0.699 & 0.047 & $+0.652$ \\
Arithmetic & 0.619 & 0.131 & $+0.488$ \\
TSV & 0.692 & 0.153 & $+0.539$ \\
TIES & 0.543 & 0.147 & $+0.396$ \\
DARE & 0.684 & 0.236 & $+0.448$ \\
\midrule
\textbf{Average} & \textbf{0.648} & 0.143 & $+0.505$ \\
\bottomrule
\end{tabular}
\end{table}

\paragraph{Analysis.}

\textbf{Pearson correlation substantially outperforms MSE.} Across both architectures and all merging methods, Pearson-based optimization vastly outperforms MSE. For ViT-B-16, the average validation correlation is 0.544 with Pearson versus 0.147 with MSE---a 3.7$\times$ improvement. For ViT-B-32, the gap is even larger: 0.648 versus 0.143, a 4.5$\times$ improvement.

This performance gap arises from the fundamental mismatch between MSE and our prediction goal. Our objective is to \emph{rank} task pairs by mergeability---predicting whether pair A will merge better than pair B---rather than predicting exact merge performance values. Correlation-based objectives directly optimize for ranking quality: they are invariant to the scale and shift of predictions, focusing entirely on whether the relative ordering is preserved.

In contrast, MSE optimization must simultaneously learn (1) the correct ranking and (2) the correct output scale. Since merge performance values are bounded between 0 and 1 (normalized accuracy), while our linear predictions $\hat{y} = \mathbf{w}^\top \mathbf{x}$ can take arbitrary values, the optimization must implicitly learn to constrain outputs to a valid range. This additional burden degrades ranking performance.

\paragraph{Conclusion.}
We adopt Pearson correlation with L1 regularization as our default objective. The correlation-based objective aligns with our ranking-focused prediction goal, while L1 regularization provides robustness against overfitting and produces interpretable sparse coefficient vectors.
